\documentclass{situationist-preprint}

\title{The Self-Undermining Argument Expanded}
\msid{MND-13-b}
\msfamily{MND-13}
\msdiscipline{Philosophy of Mind}
\mstype{Companion manuscript}
\mskeywords{eliminative materialism; belief; assertion; self-refutation; question-begging}

\begin{document}
\maketitle

\begin{abstract}
The argument of MND-13-a is that eliminativism is self-undermining: a sincere assertion of it expresses a belief, so asserting it makes it false. The standard reply is that the argument begs the question by presupposing the folk-psychological account of assertion. We develop the answer to this reply. We show that any successor notion that lets the eliminativist's utterance function as an assertion of her thesis must have the functional profile of belief, and that the eliminativist's favorite analogy, the vitalist who argues that the anti-vitalist must be alive, tells against her once it is completed. We distinguish the conclusion from the stronger claim that eliminativism is unthinkable, which we do not defend.
\end{abstract}

\section{The objection}

The argument of MND-13-a has two premises: whoever sincerely asserts a content believes it (Assumption~2.1), and someone has sincerely asserted that there are no beliefs (Assumption~2.2). Theorem~3.1 follows: there are beliefs, and eliminativism is false. The eliminativist's reply is that Assumption~2.1 belongs to the theory under dispute \citep{Churchland1981}. If folk psychology is false, its account of assertion is false too, and the argument assumes what it must prove. P.~S.~Churchland's analogy makes the point vivid: a vitalist might argue that the anti-vitalist, in asserting her thesis, shows that she is alive, hence possessed of vital spirit, hence wrong \citep{ChurchlandPS1986}. The vitalist's argument begs the question, since it assumes that being alive requires vital spirit; so, it is said, does ours.

\section{The successor dilemma}

Let $S$ be whatever state the eliminativist puts in place of belief, the state her successor theory says she is in when she sincerely utters \enquote{there are no beliefs}. We ask what $S$ must be like for the utterance to do what she needs it to do.

\begin{definition}[Belief role]\label{d:role}
A state has the \emph{belief role} with respect to a content $p$ if (i) it is correct when $p$ is true and incorrect otherwise; (ii) it is formed and revised in response to what the subject takes as evidence for $p$; and (iii) it disposes the subject to act and infer as if $p$.
\end{definition}

\begin{proposition}\label{p:dilemma}
If the eliminativist's utterance is to give her audience a reason to accept that there are no beliefs, the state $S$ it expresses has the belief role with respect to that content.
\end{proposition}

\begin{proof}
An utterance gives an audience a reason to accept a content only if the speaker's producing it is evidence for the content, and it is evidence only if it is produced by a state whose correctness condition is the content and that is responsive to evidence for it; otherwise the utterance covaries with the content no better than chance. This yields clauses (i) and (ii). The eliminativist also argues: she draws consequences from her thesis, rejects claims inconsistent with it, and designs research programs on its basis. This is clause (iii).
\end{proof}

The eliminativist must therefore take one of two options. If $S$ has the belief role, $S$ is a belief: a belief is not individuated by the vocabulary of folk theory but by the role, and a state with the role is what the word \enquote{belief} refers to on any functional account of reference \citep{Boghossian1990}. Then Assumption~2.1 holds of $S$ and Theorem~3.1 of MND-13-a goes through with $S$ in place of belief. If $S$ lacks the belief role, her utterance gives no one a reason to accept it, and her argument is not an argument; there is nothing for us to answer. Either way, the reply does not show that the original argument begs the question. It shows that the premise can be stated in neutral terms, by Definition~\ref{d:role}, which the eliminativist must accept to be engaged in assertion at all \citep{Baker1987}.

\section{Completing the analogy}

The vitalist's argument begs the question because there is an account of life that does without vital spirit: being alive is metabolizing, reproducing and maintaining a boundary. The anti-vitalist can grant that she is alive and deny the theory of life that requires vital spirit. What she eliminates is a theory of what life is, not life.

The analogy, completed, has the same structure for belief. The eliminativist may deny a theory of what belief is: that beliefs are sentence-like, that they are discrete, that their contents are those folk ascriptions assign. But the counterpart of life, in the analogy, is the role of Definition~\ref{d:role}, and she must grant that she has a state that plays it, as the anti-vitalist grants that she is alive. What she then eliminates is a theory of belief, and the conclusion that there are no beliefs does not follow from the falsity of that theory any more than the conclusion that nothing is alive follows from the falsity of vitalism. \citet{Stich1983} drew the inference from the falsity of folk theory to the nonexistence of its posits; the analogy shows why that inference fails when the posit is fixed by a role the theory's critic must herself occupy.

\begin{objection}[Pragmatic, not logical]
The argument at most shows that eliminativism cannot be coherently asserted, not that it is false. Unassertable truths are possible.
\end{objection}

\begin{reply}
That is right in general, and it is why Assumption~2.2 is needed. The argument does not infer the falsity of eliminativism from its unassertability. It infers it from the fact that it has been sincerely asserted, an event that by Assumption~2.1 involves a belief. Had no one asserted or believed the thesis, the argument would show only that it cannot be asserted truly.
\end{reply}

\begin{objection}[Unthinkability]
Then the argument proves too much: it would show that no one could even entertain eliminativism.
\end{objection}

\begin{reply}
It does not. Entertaining, supposing and considering a content do not require believing it, and nothing in MND-13-a or here denies that eliminativism can be entertained or investigated. What it cannot be is sincerely asserted truly, and since it has been, it is false.
\end{reply}

\section{Conclusion}

The charge of question-begging fails. The premise that sincere assertion expresses belief can be stated in terms of a role, and the eliminativist must occupy that role to assert anything. The vitalism analogy, completed, shows that what can be eliminated is a theory of belief, not belief. Beliefs exist, and eliminativism, having been asserted, is false.

\begin{thebibliography}{9}\small
\bibitem[Baker(1987)]{Baker1987} Baker, L.~R. (1987). \emph{Saving Belief: A Critique of Physicalism}. Princeton University Press.
\bibitem[Boghossian(1990)]{Boghossian1990} Boghossian, P.~A. (1990). The status of content. \emph{The Philosophical Review}, 99(2), 157--184.
\bibitem[Churchland(1981)]{Churchland1981} Churchland, P.~M. (1981). Eliminative materialism and the propositional attitudes. \emph{The Journal of Philosophy}, 78(2), 67--90.
\bibitem[Churchland(1986)]{ChurchlandPS1986} Churchland, P.~S. (1986). \emph{Neurophilosophy: Toward a Unified Science of the Mind-Brain}. MIT Press.
\bibitem[Stich(1983)]{Stich1983} Stich, S.~P. (1983). \emph{From Folk Psychology to Cognitive Science: The Case Against Belief}. MIT Press.
\end{thebibliography}

\end{document}
